\documentclass[12pt]{article}
\usepackage{pgf, tikz}
\usepackage{amsmath, amsfonts, amssymb, graphicx}
\usepackage{subfig}
\usepackage{float}
\usepackage[utf8]{inputenc}
\usepackage[spanish]{babel}
\usepackage{amsthm}

\setlength{\textheight}{23cm} \setlength{\evensidemargin}{0cm}
\setlength{\oddsidemargin}{-.5cm} \setlength{\topmargin}{-3cm}
\setlength{\textwidth}{17.5cm} \setlength{\parskip}{.2cm}


%opening

\begin{document}
	\begin{picture}(80, 80)
	\put(170,0){\hbox{\includegraphics[scale=0.6]{cimat_logo.png}}}
	\end{picture}
	
	\begin{center}
		\begin{huge}
			Centro de Investigación en Matemáticas, A.C.
		\end{huge}
	\end{center}

	\begin{center}
		\begin{large}
			Descripción tarea - Métodos numéricos
		\end{large}
	\end{center}
	
	\begin{center}
		\textbf{Erick Salvador Alvarez Valencia}
	\end{center}

	\begin{center}
		5 de Octubre de 2017
	\end{center}



%\maketitle

%\tableofcontents

\section{Método de potencia inversa}
El primer método a describir es el de la potencia inversa, el cuál nos ayuda a encontrar el eigenvalor más pequeño, aunque con una pequeña modificación puede encontrar algún valor que sea muy cercano a un $\delta$ dado.\\
Sea $A \in R^{nxn}$ una matriz no necesariamente simétrica, estamos interesados en encontrar un eigenpar $(\phi, \lambda)$, aplicando el método de la potencia inversa encontraremos un eigenpar tal que el valor de $\lambda$ sea el más pequeño que existe. Ahora definimos a $\delta$ como un valor que está dentro del rango $[-d, d]$ donde $d = ||A||_{\infty}$. Podemos ver entonces que si sumamos $\delta$ a la matriz $A$ obtenemos:\\
$$(A + \delta I)\phi = A\phi + \delta \phi = \lambda \phi + \delta \phi = (\lambda + \delta)\phi$$

De lo anterior se obtiene un nuevo eigenvalor asociado al mismo eigenvector. Por lo que podemos concluir que eligiendo valores de $\delta$ apropiados, podremos encontrar posiblemente todos los eigenpares de la matriz mediante el método de la potencia.\\
En el código proporcionado se realiza una discretización uniforme dentro del rango $[-d, d]$ con un tamaño de paso $\Delta d$ donde dicha constante la obtenemos como $\Delta d = 2d / N$, el usuario brindará el valor de $N$ al programa. Ahora por cada valor $\delta_i$ y con una cierta tolerancia $\tau$ llamamos a la función de la potencia, la cual nos devolverá un $\mu_i$, si se cumple que $|\mu_0 - \mu| > 0.0001$ entonces podemos considerar a $\mu$ y a su respectivo $\phi$ como un eigenpar válido. Para asegurarnos de tener una muy buena probabilidad de encontrar todos los eigenpares de una matriz, ejecutamos el método de la potencia con un máximo de 1000 iteraciones.

\subsection{Ejemplo de ejecución}
A continuación se mostrará el resultado de la ejecución del método de la potencia con una matriz de $5x5$.



\subsection{Compilación y ejecución}
\textbf{Para compilar:} En la carpeta encontraremos los archivos $.c$ y $.h$ con los que se podrá compilar el ejecutable. De la misma forma, en conjunto con los archivos anteriores, también podremos encontrar un Makefile para, en caso de encontrarse en linux, compilar de manera sencilla.

\begin{enumerate}
	\item \textbf{Compilar usando Makefile:} En la terminal, nos colocamos en el directorio donde se encuentre el programa, y ejecutamos el comando $make$, automáticamente se realizará la compilación y se generará el ejecutable.\\
	\item \textbf{Compilar directamente:} De la misma forma, podemos compilar directamente usando los siguientes comandos (en terminal):\\
	\begin{itemize}
		\item gcc -c main.c
		\item gcc -c memo.c
		\item gcc -c matriztridiagonal.c
		\item gcc -c reader.c
		\item gcc -o main main.o memo.o matriztridiagonal.o reader.o -lm
	\end{itemize}
\end{enumerate}

\textbf{Para ejecutar:} Únicamente debemos de usar el comando $./main$ para ejecutar el programa en consola.\\
El programa recibe un argumento, un entero, el número de particiones que se realizarán. A continuación se mostrará un ejemplo de cómo ejecutar el programa:\\

$./main\ 10$

\section{Ejemplo de ejecución}

A continuación se mostrarán ejemplos de ejecución del programa con 10 y 100 particiones:\\

\end{document}
